% !TeX root = ../main.tex
\chapter{Background and Related Work}

\section{Foundations of Neural Chess Evaluation}

\subsection{The Role of Evaluation and Classical Approaches}
At the core of every chess engine lies an evaluation function $f: \mathcal{S} \rightarrow \mathbb{R}$ that assigns a scalar value to a board position $s \in \mathcal{S}$, indicating the expected outcome from that position (typically from the perspective of the player to move). This value guides the search algorithm -- usually a variant of \textbf{alpha-beta search} with iterative deepening -- by pruning unpromising branches and selecting the most best moves.

For decades, chess evaluation functions were handcrafted. A classical evaluator is a weighted linear combination of chess-specific features:
\[
f(s) = \sum_{k} w_k \cdot \phi_k(s)
\]
where $\phi_k(s)$ are feature functions and $w_k$ are scalar weights. Typical features include material balance, piece-square tables (PSTs), mobility, king safety, and pawn structure. Chess experts carefully tuned these features over decades, using a combination of intuition, empirical testing, and later automated optimization techniques.

Traditional engines typically evaluate a position using a scalar value called \emph{centipawn}. Positive values indicate an advantage for White, while negative values indicate an advantage for Black. The name hints at the fact that the value of a pawn corresponds to 100 centipawns, but this is not entirely accurate, as the value of each piece changes dynamically depending on the position. Therefore, interpretability is a crucial weakness of this metric. Similarly, since centipawns hint at material advantage, they lose meaning when the probability of a color winning is extreme, usually capping at around 5000 centipawns.

% todo reference to some centipawn discussion online?

A notable example of a highly tuned classical evaluator is \textbf{PeSTO} (Piece-Square Tables Only), an evaluation function by Ronald Friederich that relies exclusively on piece-square tables. Its tables are optimized via \emph{Texel}'s tuning method and use a \emph{tapered evaluation} that interpolates between separate opening and endgame tables based on the game stage. Conceptually, PeSTO is equivalent to a linear feed-forward network: both compute a weighted sum of piece-square features, with no interaction terms between pieces. While still a handcrafted linear form, PeSTO demonstrates how far such classical approaches can be pushed through data-driven tuning \parencite{friederichpesto}.

\begin{figure}[H]
	\centering
	\includegraphics[width=0.86\textwidth]{../thesis-plots/pesto_pawn_mg.png}
	\caption{PeSTO middlegame pawn table, in centipawns.}
	\label{fig:pesto-pawn-mg-1}
\end{figure}

While classical evaluators are extremely fast, they suffer from fundamental limitations. Human-designed features encode human assumptions about chess, which may not align with the optimal understanding of the game. Moreover, these models do not take into account the positions as a whole; the same parameters give a positional bonus for a central knight whether the position is a quiet middlegame or a tactical melee, even though the knight's practical value may differ dramatically across phases.

Classical evaluators are also vulnerable to the \emph{horizon effect}. A search algorithm constrained to a fixed depth evaluates positions at the search frontier as if the game were stable at that point. If a tactical sequence begins at the horizon but extends one ply further, the engine cannot see the consequences and may assign an inaccurate evaluation. Classical evaluators, which rely on static features, are particularly susceptible: they treat a position at the horizon as if it were quiet, ignoring the tactical possibilities that would unfold if the search continued.

\subsection{The Shift to Neural Evaluation}
The limitations of handcrafted features led to the adoption of neural networks as evaluation functions. \textbf{AlphaZero} \parencite{silver2018} demonstrated that a deep convolutional network, trained via self-play reinforcement learning, could surpass the best classical engines. However, these networks are computationally expensive -- requiring millions of operations per evaluation -- making them unsuitable for resource-constrained devices or for engines that must evaluate millions of positions per second.

AlphaZero's value head outputs a scalar $v\in[-1,+1]$, interpreted as the expected game outcome from the current player's perspective. This scalar is the training target for the value network. However, AlphaZero's training and inference rely on \textbf{Monte Carlo Tree Search (MCTS)}, which builds a search tree by repeatedly simulating trajectories and using the neural network to evaluate leaf nodes. This process requires many forward passes of the network per position, making it computationally expensive and poorly suited for engines that evaluate millions of positions per second with alpha-beta search. The AlphaZero network itself is a deep residual architecture with millions of parameters, requiring floating-point operations and substantial memory, which far exceeds the capacity of microcontrollers \parencite{browne2012}.

% todo reference to MCTS?

\begin{figure}[H]
	\centering
	\includegraphics[width=1.0\textwidth]{../thesis-plots/mcts_wiki.png}
	\caption{The four steps of Monte Carlo tree search.}
	\label{fig:mcts-wiki-2}
\end{figure}

\subsection{NNUE: A Hybrid Approach}
To bridge the gap between the representational power of neural networks and the speed of classical evaluators, the \textbf{Efficiently Updatable Neural Network (NNUE)} architecture was introduced in the \emph{shogi} engine \emph{YaneuraOu} and later adopted by \emph{Stockfish}. 

% todo one phrase for shogi?

The architecture consists of two ingredients: a sparse first layer called the \emph{accumulator}, and a small fully-connected head. The accumulator layer maps a binary representation of the board to a hidden representation $h \in \mathbb{R}^{d}$ via $h = W_{L1} \cdot x$, where $x$ is a sparse binary vector indicating the presence of pieces on squares. The key insight is that a chess move changes only a few bits of $x$, so $h$ can be updated incrementally by adding and subtracting the corresponding columns of $W_{L1}$, rather than by recomputing the entire matrix-vector product from scratch. This makes the cost of the computation of L1 negligible.

The second ingredient, the fully connected head, is typically one or two hidden layers followed by an output layer, mapping $h$ to the position value. A complete NNUE engine leverages alpha-beta search with iterative deepening: starting from the root position, the engine searches deeper and deeper, using the NNUE evaluation at leaf nodes to guide the pruning and ordering of moves. At every node of this search tree, the evaluation function is called hundreds or thousands of times, making its speed critical. Being trained on enormous amounts of data, a NNUE does not suffer from the horizon problem as much as a PST does.

Beyond the standard NNUE formulation, recent trends in machine learning suggest that the handcrafted sparse accumulator could be replaced by a learned \emph{world model} encoder. Instead of relying on predefined piece-square features, a feedforward neural network could process the raw board state to learn a continuous, adaptable spatial representation. While this approach sacrifices the incremental update efficiency of the traditional accumulator, it shifts the learning effort entirely to a shared representation, potentially offering a more flexible foundation for complex positional understanding.

% todo support claim "recent trends in machine learning suggest that..."

\subsection{Value Targets: Centipawns, Expected Value, and WDL}
A fundamental design choice in any chess evaluation function is the representation of the target value. Three main approaches have emerged in practice: scalar centipawns, scalar expected value (EV), and full WDL distributions.

Classical evaluators and Stockfish-style NNUEs output a scalar value in centipawns. This representation is simple and interpretable, but it compresses the uncertainty of the game outcome into a single number: a position with a 100 centipawn advantage might be a quiet positional advantage, or a tactical trap, both of which require different handling during search. Moreover, training a network to predict centipawns typically uses mean squared error, which treats all deviations equally regardless of whether they lie near the decision boundary.

AlphaZero adopted a different scalar representation: $v \in [-1, +1]$, interpreted as the expected game outcome. This value is the difference between the probability of a win and the probability of a loss: $v = p_W - p_L$. This representation is more naturally calibrated to game outcomes than centipawns and avoids arbitrary scaling. However, like centipawns, it compresses the full distribution into a single scalar, losing information about the probability of a draw. A position with $v = 0$ could be a drawish endgame, or a middlegame position where the engine is equally uncertain about win and loss.

The approach adopted in this work is to train the NNUE to output a full probability distribution over the three possible game outcomes: Win, Draw, and Loss. The output head produces three logits, converted to probabilities via softmax:
\[
p_{\text{WDL}}(s) = (p_W, p_D, p_L), \qquad p_W + p_D + p_L = 1
\]
The scalar expected value used during search is then derived as $v = p_W - p_L$, but the network is trained to match the full distribution rather than just the scalar. This representation has several advantages:
\begin{itemize}
	\item It preserves information about uncertainty: a position with $p_W = 0, p_D = 1, p_L = 0$ has the same scalar value as one with $p_W = 0.5, p_D = 0, p_L = 0.5$, but the two positions are fundamentally different.
	\item It enables the use of \textbf{soft cross-entropy} as the loss function, which provides a richer training signal than mean squared error: the network is encouraged to match the full shape of the distribution, not just its mean.
	\item The loss does not saturate at extreme values, as squared error on $v$ would. Finally, the softmax output is naturally calibrated as a probability distribution, avoiding the need for an extra computational layer to derive EV from raw outputs.
\end{itemize}

\subsection{Why Evaluation Must Be Cheap}
The search tree explored by an alpha-beta engine grows exponentially with depth. Even with effective move ordering and pruning techniques, the number of evaluated positions increases dramatically as the search goes deeper. A chess engine that searches deeper consistently outperforms one that searches shallower, as each additional ply reveals tactical patterns and strategic nuances that would otherwise remain hidden.

\begin{figure}[H]
	\centering
	\includegraphics[width=0.86\textwidth]{../thesis-plots/alpha_beta_search_ChatGPT_Image_Jun_1_2026_01_01_16_PMQZ7DEM.png}
	\caption{Alpha-beta search and pruning.}
	\label{fig:alpha-beta-search-chatgpt-image-jun-1-2026-01-01-16-pmqz7dem-4}
\end{figure}

This means that any overhead added to the evaluation function is directly subtracted from the engine's effective search depth. If the evaluator becomes slower, the engine must reduce its search depth to maintain the same response time -- sacrificing playing strength in the process.

On an embedded device, this constraint is even more severe. Memory is limited, floating-point operations are expensive, and every instruction counts. Any routing mechanism for a mixture-of-experts NNUE must add only a trivial cost -- ideally, a simple matrix-vector product or a table look-up -- to avoid degrading the engine's search performance. In short, the evaluation function must be expressive enough to assess positions accurately, yet cheap enough to be called millions of times during a game. This trade-off is the central engineering challenge addressed by this thesis.

\subsection{Current State of the Art for Tiny Hardware: Cfish}
For resource-constrained devices, the most relevant reference implementation is \textbf{Cfish}, a port of the Stockfish chess engine written in plain C by Ronald de Man. While Stockfish is written in C++ and targets desktop-class hardware, Cfish strips away the C++ abstractions and compiles to a much smaller binary, making it viable for microcontrollers and other embedded platforms. The engine can be compiled with various evaluation backends, including a pure NNUE mode that excludes the classical handcrafted evaluation entirely, resulting in an even smaller executable \parencite{deman2016}.

From an architectural standpoint, Cfish shares the same core search and evaluation logic as Stockfish. The input to the evaluation function is a board position represented internally as a compact bitboard structure. The NNUE evaluation itself follows the \emph{halfkp} feature set: for each piece on the board, the network considers the piece's square together with the square of the king of the same color, producing a set of activated features that are fed into the accumulator. The accumulator is updated incrementally as moves are made, avoiding recomputation from scratch.

Cfish is particularly relevant to this thesis for two reasons. First, it demonstrates that a full-featured NNUE engine \emph{can} be made to run on constrained hardware, provided the implementation is careful about memory layout and avoids C++ overhead. Second, it serves as a practical baseline: any Mixture-of-Experts NNUE architecture proposed for tiny devices should be at least as fast and compact as Cfish in its pure NNUE mode, while offering improved evaluation accuracy through expert specialization.

\section{Mixture of Experts and Specialization Paradigms}

The \textbf{Mixture of Experts (MoE)} architecture is a neural network design pattern in which multiple specialized sub-networks, or \emph{experts}, are combined through a routing mechanism that selects or weights their contributions based on the input. The central idea is that distinct regions of the input space may require different processing, and dedicating separate capacity to each region can improve overall performance without substantially increasing the cost of a single forward pass.

\begin{figure}[H]
	\centering
	\includegraphics[width=0.86\textwidth]{../thesis-plots/nnue_architecture_moe.png}
	\caption{NNUE with a mixture-of-experts head.}
	\label{fig:nnue-architecture-moe-5}
\end{figure}

\subsection{Fixed vs. Learned Routing}
MoE architectures can be broadly divided into two categories based on how the routing is determined: \emph{fixed routing} and \emph{learned routing}.

In fixed routing schemes, the assignment of inputs to experts is determined by a predefined rule, often based on domain knowledge. In chess engines, this corresponds to handcrafted bucketing: positions are assigned to buckets based on material count, piece presence, or game phase. The routing is deterministic, interpretable, and computationally cheap, but it relies on human assumptions about which regions of the state space are meaningfully distinct. The rule is fixed after design and cannot adapt to the data.

In learned routing schemes, a trainable \emph{gating network} (or router) learns to assign inputs to experts based on the input features themselves. The gating network typically produces a probability distribution over experts, and the final output is a weighted combination of expert outputs, or a single expert selected by argmax. This approach is more flexible: the router can learn to assign inputs to experts in ways that may not align with human intuition, potentially discovering structure in the data that handcrafted rules would miss. However, learned routing introduces additional parameters and computational cost, and it may require careful design to avoid load imbalance or mode collapse.

% todo quote "weighted combination of expert outputs"

\subsection{Applications in Vision, NLP, and Reinforcement Learning}
Mixture of Experts has been successfully applied across a wide range of domains. In natural language processing, large-scale MoE models such as the Switch Transformer \parencite{fedus2021} and GLaM \parencite{du2022} achieve state-of-the-art performance by scaling the number of parameters while keeping per-token computation constant: only a subset of experts is activated for each token. In computer vision, MoE architectures have been used to efficiently scale convolutional networks and vision transformers, where experts specialise in various visual patterns or object classes. In reinforcement learning, MoE has been applied to multi-task and multi-domain settings, where the experts specialise in different tasks or environments, and a gating network selects the appropriate expert for the current context.

Despite their success in these domains, MoE architectures are not very often deployed in resource-constrained settings such as microcontrollers. The gating network and the multiple expert heads introduce memory and computational overhead that is acceptable on servers but prohibitive on embedded devices. Moreover, many MoE implementations rely on sparse activation and require specialised hardware support for efficient execution. These constraints are less severe in the chess domain, where the evaluation function must be simple and fast, but they still shape the design choices of this work.

\subsection{Parameter-Efficient Specialization: From LoRA to NNUE Heads}
A fundamental challenge in Mixture-of-Experts architectures is how to allocate specialized capacity without incurring the massive memory and computational costs of duplicating the entire network. The solution, widely adopted in modern deep learning, is to freeze a large, shared representation and restrict the specialized parameters to a small, adaptable subset.

In the context of large language models, this paradigm is best exemplified by \textbf{Low-Rank Adaptation (LoRA)} \parencite{hu2021}. LoRA freezes the base model's weights and injects trainable low-rank matrices into specific layers, enabling efficient fine-tuning with a tiny fraction of the original parameters. Recent works have successfully combined this parameter-efficient specialization with gradient-informed routing. \textbf{ELREA} \parencite{li2025} partitions training data by gradient directions and fine-tunes a separate LoRA expert on each partition, while \textbf{GradientSpace} \parencite{sridharan2025} clusters LoRA gradients and uses a lightweight router to select the appropriate adapter at inference. These approaches demonstrate that isolating the learning signal into a small set of adaptable parameters is a highly effective way to manage optimization conflicts and enable expert specialization.

Although LoRA is directly applicable to the resource-constrained NNUE setting, the exact mathematical mechanism is tailored for the massive transformer layers of LLMs. In our architecture, the shared representation (the L1 accumulator) is already the most memory-intensive component, and the specialized evaluation head comprises only a few thousand parameters. Applying low-rank adapters to such a small head would introduce unnecessary computational overhead and complicate the quantization required for microcontroller deployment.

Instead of low-rank matrices, this thesis adopts the core structural principle of LoRA: freezing the shared trunk and duplicating only the small, task-specific head. Each expert is implemented as a separate, full-rank instance of the L2 and output layers, initialized from the base head and fine-tuned on its assigned bucket. By keeping the L1 accumulator frozen and shared, the architecture preserves the exact benefits of the LoRA paradigm -- massive parameter efficiency, isolated specialization, and a shared foundational representation -- while remaining strictly compatible with the extreme memory and latency constraints of embedded chess engines.

\subsection{Adapting MoE to Embedded Chess}
From the broader MoE literature, the key ideas that inform this work are: the principle of partitioning the input space to enable specialisation; the distinction between fixed and learned routing; and the observation that gradient-based clustering can be used to discover meaningful partitions. However, the specific constraints of embedded chess engines impose a different set of trade-offs. The routing mechanism must be relatively inexpensive, ruling out large gating networks; the expert heads must be small and easily quantizable; and the partition must be learned from value estimation targets rather than language model training data. This thesis adapts the gradient-based partitioning paradigm to these constraints, proposing a lightweight linear dispatcher that routes positions to specialised heads with negligible overhead.

\section{Sample Gradients as a Learning Signal}

\subsection{Per-Sample Gradients}
In standard neural network training, the loss is computed over a mini-batch, and the resulting gradient is an average over the examples in that batch. This averaging discards information about how individual examples differ in their contribution to the update, treating the batch as a single unit. Per-sample gradient computation, by contrast, produces a separate gradient vector for each example in the batch, capturing the direction in which the model parameters would need to move to reduce the loss on that specific example alone.

Per-sample gradients have found applications in differential privacy, where individual example contributions must be bounded to protect privacy, and in meta-learning, where the gradient with respect to a specific example informs task adaptation. More recently, a line of work has used per-sample gradients as a representational tool: the gradient of a single example encodes what that example ``wants'' the model to learn, and clustering these gradients reveals groups of examples that share similar learning requirements.

A key property of per-sample gradients is that they are directly defined in parameter space. The gradient 
\[
\Delta_i=\nabla_w \, \mathcal{L}(f_w(s_i),v_i)
\]
specifies a direction in the same space in which the model's weights reside. This distinguishes them from hidden-layer activations, which characterize the representation that the model currently produces for an input but say nothing about how that representation should change. Two examples may produce very similar activations because the model already evaluates them well, yet require different updates to correct subtle errors; conversely, examples with dissimilar activations may benefit from similar parameter adjustments. The gradient, by construction, resolves this distinction: it is the direction in which the loss would decrease most rapidly for that example, and therefore a direct measure of the learning signal.

The connection between per-sample gradients and expert specialization is supported by both theoretical and empirical work. \textcite{kawata2025} prove that a Mixture of Experts trained with stochastic gradient descent can provably detect and exploit latent cluster structure in the data, dividing the problem into easier subproblems and achieving sample complexity gains over a single model. This result provides formal justification for using gradient information to guide the partitioning of a state space. Empirically, ELREA and GradientSpace demonstrate that clustering gradients yields partitions that support effective expert specialization in language models. TAGCOS extends the idea to coreset selection, using gradient clustering to identify representative examples for efficient fine-tuning.

\subsection{Gradient Similarity and Specialization}
If per-sample gradients encode what each example wants the model to learn, the natural next question is what it means for two gradients to be similar, and why similarity should matter for specialization. The answer lies in the geometry of the gradient space and in the optimization dynamics it induces.

Two per-sample gradients are similar when they point in approximately the same direction in parameter space, that is, when their cosine similarity is high. This means that the parameter updates that would reduce the loss on one example would also reduce the loss on the other. When many examples share this property, they can be served by the same parameter configuration without conflict: a single update direction benefits all of them simultaneously. When gradients point in opposing directions, the update demanded by one example harms the other, and a single model must compromise, settling on parameters that are suboptimal for both.

Clustering per-sample gradients groups together examples whose learning requirements are aligned, so that each cluster can be served by a dedicated expert without internal conflict. The partition is not defined by what the examples look like in input space, but by what they demand of the model in parameter space. The connection to expert specialization is made explicit in the Mixture of Experts framework: if two examples are assigned to the same expert and their gradients conflict, the expert will be forced into a compromise; if they align, the expert can specialize without tension. Maximizing the separation between clusters in gradient space minimizes the conflict within each expert and maximizes the distinctness of what each expert learns.

\subsection{Gradient-Based Clustering}
Gradient-based clustering answers the need for aligned learning requirements by treating the per-sample gradient as a data representation and applying standard clustering algorithms to the resulting set of vectors. Each example is mapped to a point in the gradient space, and clusters correspond to groups of examples whose gradients point in similar directions.

The appeal of this approach lies in its directness. Unlike clustering hidden-layer activations, which groups inputs by what the model already represents, gradient clustering groups examples by what the model needs to change. The partition is therefore aligned with the optimization objective rather than with the current state of the representation.

A practical challenge arises from the dimensionality of the gradient space. The gradient vectors for a single example are as large as the parameter vector itself, and storing them for millions of examples is a non-trivial computational burden. Two strategies are commonly used to address this. The first is to reduce the dimensionality of the gradients before clustering, for example through random projections or principal component analysis. ELREA adopts random projection to map gradients to a lower-dimensional space before clustering. GradientSpace instead operates directly on the full-dimensional gradient space, using an online SVD-based algorithm to identify latent skills without materializing all sample gradients at once. 

The second strategy, adopted in this work, is to normalize the gradients before clustering, projecting each vector onto the unit sphere. Normalization removes magnitude information and focuses the clustering on direction, which is the signal most relevant to specialization. This aligns with the observation that the direction of the update determines whether two examples are compatible, while the magnitude is sensitive to the current loss value and position difficulty. Normalization also makes the clustering more robust to outliers and aligns with the cosine-distance objective used in the diversity metric.

% todo maybe summarize "Gradient Similarity and Specialization" and "Gradient-Based Clustering", and make them more distinct?

\section{State-Space Bucketing in Chess}

\subsection{Handcrafted Bucketing}
In NNUE-based chess engines, the state space is often partitioned into discrete buckets using manually designed rules. Each bucket corresponds to a region of the state space where positions are assumed to share similar characteristics, and a separate set of output weights is trained for each bucket. This approach, sometimes referred to as \emph{phase-based bucketing}, is a pragmatic compromise: it allows the evaluation function to adapt to various types of positions while keeping the per-head network small and fast.

\begin{figure}[H]
	\centering
	\includegraphics[width=0.86\textwidth]{../thesis-plots/opening_middlegame_endgame.png}
	\caption{Opening, middlegame, and endgame positions.}
	\label{fig:opening-middlegame-endgame-6}
\end{figure}

Common features used for bucketing include the total material count, the presence or absence of specific pieces such as queens or bishops, the location of the kings, and the overall game phase derived from material. For example, Stockfish historically used a phase classification that interpolates between opening and endgame evaluations based on material remaining. The Kaggle FIDE \& Google Efficient Chess AI Challenge popularised a variant where the game is divided into three phases based on material thresholds, with separate evaluation heads for each phase \parencite{kagglefide}.

% todo reference to youtube video https://youtu.be/Ne40a5LkK6A

These handcrafted bucketing schemes are computationally cheap and interpretable. However, they suffer from fundamental limitations. The features are designed based on domain expertise, which may not align with the actual structure of the learning signal. A position in the middlegame with a queen on the board may require a very different adjustment to the evaluation head depending on whether it is a quiet positional struggle or a tactical melee -- yet both are assigned to the same bucket based on material count. Conversely, two positions that appear superficially distinct may require similar adjustments to the head, but are placed in different buckets. In other words, handcrafted features encode what a programmer assumes are distinct regimes, not what the model actually needs in order to specialise.

\subsection{The Need for Data-Driven Bucketing}
The limitations of handcrafted bucketing have motivated research into data-driven alternatives, where the partition of the state space is learned directly from the training data rather than prescribed by human assumptions. These approaches can be broadly divided into two categories, both of which present significant challenges for embedded chess engines.

One natural approach is to cluster the activations of the network's hidden layers, grouping positions that the model already represents similarly. The intuition is that positions that produce similar hidden representations are likely to require similar processing from the subsequent layers. This approach has the advantage of being fully data-driven and requiring no chess-specific feature engineering. However, it suffers from a fundamental limitation: activations capture what the model \emph{knows}, not what it \emph{needs} to adjust. Two positions may produce similar hidden representations because the model already evaluates them correctly, yet require very different updates to the head; conversely, positions with dissimilar activations may require similar adjustments. The partition is defined by the model's current state, not by the learning signal that drives specialisation.

% todo check how often we say this thing about "knows" vs "need" -- repeating a couple times is OK

An alternative to clustering is to learn a routing policy that directly selects among experts based on the input state. In the broader machine learning literature, mixture-of-experts architectures typically employ a learnable gating network. In the chess domain, recent work has explored learned routing for expert selection. \textbf{Hexa\"issa} \parencite{ngo2026} formulates expert selection as a MoE problem, learning a gating policy that dynamically selects among heterogeneous state-of-the-art engines. \textbf{M2CTS} \parencite{helfenstein2024} combines MoE with MCTS, using a modular framework that adapts strategy dynamically based on game phase. These approaches demonstrate that learned routing can be effective in chess, but they typically operate at the level of entire engines or large networks. Moreover, they rely on expensive routing mechanisms that would be prohibitive at every node of an alpha-beta search.

This creates a clear gap: existing learned bucketing methods either capture the wrong signal (activations) or are computationally prohibitive (gating networks). This thesis addresses this gap by proposing a method that learns the partition from sample gradients -- the learning signal itself -- and then distills this partition into a lightweight linear dispatcher that adds negligible overhead at inference time.

\section{Related Work and Positioning}

The preceding sections have established the conceptual foundations of this thesis. This final section places the proposed method within the broader body of work on efficient chess engines, teacher-student evaluation, and gradient-based expert specialization, clarifying how the method presented in this work differs from existing approaches.

\subsection{Efficient Chess Engines and Quantization}
The pursuit of efficient chess evaluation has a long history, driven by the need to evaluate millions of positions per second on commodity hardware. The introduction of NNUE marked a turning point by demonstrating that a neural network could be evaluated within the same latency budget, provided that the first layer is sparse and updated incrementally. This principle allowed Stockfish and its derivatives to adopt neural evaluation without sacrificing search speed.

The portability of NNUE to constrained environments was demonstrated by Cfish, a C port of Stockfish that strips away the C++ abstractions and compiles to a much smaller binary. Cfish provides multiple evaluation backends, including a pure NNUE mode that eliminates the classical evaluation entirely, making it viable for embedded platforms. It also incorporates SIMD-optimized (Single Instruction, Multiple Data) attack generation for AVX2 and AVX-512 (instruction sets for parallel vector processing), further reducing the cost of move generation and evaluation. Cfish is the closest reference implementation to the target of this thesis, and it serves as the host engine into which the MoE NNUE is integrated in Chapter 4.

Beyond the choice of architecture, efficiency in chess engines is achieved through a combination of search pruning, reduction techniques, and integer quantization. Modern engines employ a family of pruning methods to reduce the effective branching factor of the alpha-beta search. Quantization, by contrast, directly concerns the evaluation function: NNUE networks are designed for low-precision integer inference, using int8 and int16 operations to exploit the available hardware performance of modern CPUs. In Chapter 4, we adopt the same quantization strategy when deploying the MoE NNUE, ensuring that the additional expert heads and dispatcher fit within the memory and latency budget of the device.

What distinguishes this thesis from the efficient-engine literature is not the quest for raw speed, but the question of how to allocate a fixed parameter budget more effectively. The engines discussed above use a single head to evaluate all positions; the MoE architecture distributes the head parameters across multiple experts and adds a dispatcher to route each position to the appropriate expert. Whether the specialization gained from this partition outweighs the additional memory footprint is the central empirical question of Chapter 5.

\subsection{Teacher--Student Evaluation}
The idea of training a compact model to imitate a stronger or more expensive one has a long history in machine learning, formalised most influentially by \textbf{\textcite{hinton2015}} under the name of knowledge distillation. In the original formulation, a large teacher network produces soft probability distributions over classes, and a smaller student network is trained to match those distributions rather than the hard labels of the training data. The soft targets carry more information than one-hot labels because they encode the teacher's uncertainty and the relative similarity between classes.

In the chess domain, teacher--student evaluation has become the dominant paradigm for training evaluation functions. AlphaZero used a form of self-distillation: the same network served as both policy and value function, and the targets for training were generated by MCTS guided by the network itself. Lc0, the open-source implementation of the AlphaZero approach, follows the same principle and has produced some of the strongest chess networks available. Its value head outputs a probability distribution over win, draw, and loss, making it directly compatible with the soft cross-entropy objective used in this work.

A different teacher--student setup is used in the training of Stockfish's NNUE. The teacher is not a separate network but the engine itself, run at a shallow search depth on a large corpus of positions. The resulting centipawn evaluations, or in newer versions the WDL probabilities derived from the search, serve as targets for the NNUE. 

In this thesis, we adopt Lc0 as the teacher. First, its WDL output aligns naturally with the three-way output of the NNUE value head. Second, its strength makes it a reliable source of labels across all phases of the game. What distinguishes this thesis from the teacher--student literature is not the use of distillation itself, which is standard practice, but the way in which the student is partitioned. Rather than training a single student to match the teacher on all positions, we train multiple students, each specialised on a distinct region of the state space. The specialisation arises from the partition of the training data, not from a change in the target.

\subsection{Gradient-Informed Specialization and the Research Gap}
The works most directly related to this thesis are those that use gradient information to partition training data and train experts on the resulting groups. These approaches share a common pipeline: compute per-sample gradients, cluster them, train an expert on each cluster, and deploy a routing mechanism to select the appropriate expert at inference.

ELREA partitions language model training data by gradient direction and trains a separate LoRA expert on each partition. The motivation is to reduce optimization conflicts: when examples in a batch demand opposing parameter updates, the resulting gradient is a compromise that benefits none of them. By grouping examples with aligned gradients, ELREA ensures that each expert sees a coherent learning signal. At inference, the method routes to a weighted ensemble of experts based on gradient similarity, requiring on-the-fly gradient computation for the input. This makes inference expensive and unsuitable for latency-critical applications.

\textbf{GradientSpace} \parencite{sridharan2025} follows a similar pipeline but addresses the inference cost by training a lightweight encoder-based router. The router predicts the appropriate expert from the input features alone, eliminating the need for gradient computation at inference time. The method clusters LoRA gradients using an online SVD-based algorithm and trains a separate LoRA expert per cluster. This is the closest methodological precedent for the work presented in this thesis.

\textbf{TAGCOS} \parencite{zhang2024} applies gradient clustering to coreset selection rather than expert training. It clusters per-sample gradients and selects a representative subset from each cluster, reducing the amount of data needed for fine-tuning while preserving the diversity of the learning signal. TAGCOS is relevant because it confirms that gradient-based grouping is a robust signal for data selection.

\textbf{Gradient Atoms} \parencite{rosser2026} takes a different approach, decomposing per-document gradients into sparse components via dictionary learning in a preconditioned eigenspace. The goal is behaviour discovery and attribution rather than expert training, but the work reinforces the idea that gradients contain interpretable structure that can be exploited for downstream tasks.

Several key distinctions separate this thesis from these prior works. First, regarding the expert implementation, prior works rely on LoRA adapters for massive transformer models, whereas this work employs separate, full-rank NNUE heads sharing a frozen L1 accumulator, as the base model is already highly compact. Second, concerning the routing mechanism, while ELREA requires on-the-fly gradient computation and GradientSpace uses a small encoder, this thesis distills the partition into a lightweight linear dispatcher on frozen L1 activations, ensuring compatibility with microcontroller deployment. Third, the domain shifts from token distribution in language models to WDL value estimation in chess. Finally, this work normalizes gradient vectors to the unit sphere before clustering, focusing the partition on the direction of the update rather than its magnitude, which is sensitive to position difficulty and loss scale.

In summary, this thesis adapts the gradient-based expert specialization paradigm to a new domain and a new set of constraints. The core idea -- cluster by gradient direction, train experts on the clusters, route at inference -- is shared with ELREA and GradientSpace. The contributions are the adaptation to NNUE, the use of head-parameter gradients rather than LoRA gradients, the linear dispatcher on frozen L1 activations, and the validation on a resource-constrained chess engine. These differences reflect the distinct requirements of embedded inference, where every additional parameter and every additional operation has a measurable cost.